\documentclass[11pt]{letter}
\usepackage[utf8]{inputenc}
\usepackage[T1]{fontenc}
\usepackage[margin=1in]{geometry}
\usepackage{siunitx}
\usepackage{hyperref}
\urlstyle{same}

\signature{Jonathan Mace\\Corresponding author}
\address{Jonathan Mace\\Microsoft Research\\Redmond, WA\\jonathanmace@microsoft.com}
\date{\today}

\begin{document}

\begin{letter}{%
Editor\\
\textit{Postharvest Biology and Technology}}

\opening{Dear Editor,}
\small

Please consider our manuscript, ``Can You Hear Rind Stiffness? A Shell-Mechanics Model Linking Watermelon Tap Tone to Effective Rind Stiffness,'' for publication in \textit{Postharvest Biology and Technology}. The paper develops an equivalent-sphere shell model for watermelon tap-tone acoustics that links the measured resonance frequency to an effective rind stiffness parameter. Across the analysed parameter space, Sobol sensitivity analysis shows that rind elastic modulus is the dominant acoustic driver, with total-order index $S_T = 0.54 \pm 0.05$. In the practically important low-turgor limit, the governing expression admits a closed-form inversion from measured frequency to effective stiffness, turning a familiar field heuristic into a transparent mechanics-based estimator.

We believe the manuscript is a strong fit for \textit{Postharvest Biology and Technology} because it provides a mechanistic postharvest framework rather than another purely empirical classifier. The model explains why tap tone tracks rind mechanics, compares predictions across multiple cultivar geometries within a common analytical framework, and identifies the conditions under which inversion is well-conditioned and experimentally useful. This should interest readers working on non-destructive quality assessment, fruit firmness, postharvest mechanics, and acoustic sensing.

We also emphasise the paper's scope carefully. The model does \emph{not} claim that frequency alone is a direct ripeness measure. Instead, it establishes the mechanistic basis for future cultivar-specific calibration by showing how resonance data can be converted into an effective stiffness parameter that may then be related to maturity using destructive reference measurements. We hope this combination of analytical clarity, practical inversion, and restrained interpretation will make the manuscript valuable to the journal's readership.

The manuscript is original, has not been published previously, is not under consideration elsewhere, and has been approved by both authors. Brian R. Mace's submission email is \texttt{b.mace@auckland.ac.nz}.

\closing{Yours sincerely,}

\end{letter}
\end{document}
